%!TEX root = ../template.tex
%%%%%%%%%%%%%%%%%%%%%%%%%%%%%%%%%%%%%%%%%%%%%%%%%%%%%%%%%%%%%%%%%%%%
%% chapter4.tex
%% NOVA thesis document file
%%
%% Chapter with lots of dummy text
%%%%%%%%%%%%%%%%%%%%%%%%%%%%%%%%%%%%%%%%%%%%%%%%%%%%%%%%%%%%%%%%%%%%

\typeout{NT FILE chapter4.tex}%

\chapter{OCaml-Flat and OFLAT Systems}
\label{cha:ocaml_flat_and_oflat_systems}


The OCamlFLAT library and the OFLAT platform are the main componentsof this thesis, both developed in OCaml.
They are already in an advanced stage of development, thanks to the work of several students and professorswho have contributed to this project.

The OCamlFLAT library consists of various modules that make up the logical part of the system, where several FLAT theory models have already been implemented. The OFLAT platform, on the other hand, is a web application that provides the system’s graphical interface.


\section{OCamlFLAT}
\label{sec:ocamlflat}

The OCamlFLAT library is organized into modules to ensure a clearer and more structured organization. There are modules that implement models of FLAT theory, such as regular expressions, finite automata, and context-free grammars, as well as auxiliary modules for tests, error handling, and parsers. It is important to note that the modules follow a hierarchical relationship among themselves, which facilitates the maintenance and expansion of the library \cite{OFLAT}.

Some of the most important modules are:

Entity: This module represents the most abstract level of the library. The entities are the exercises and the FLAT models, each containing a type, a description, and a unique ID.

Exercise: This module is responsible for supporting exercises, which are validated using UnitTests. It extends the Entity module, inheriting its properties and methods.

Model: An abstract module that also extends Entity and introduces common functionalities to the FLAT models available in the library. Examples of methods included in this module are validate, which checks the validity of a model, and accept, which processes inputs according to the model.

Examples of modules that represent FLAT models and extend Model include:

RegularExpression: This module contains all the functionalities needed to work with regular expressions, allowing the definition and manipulation of these expressions.

FiniteAutomaton: This module implements the functionalities of a finite automaton, including the definition of states and transitions.

ContextFreeGrammar: This module offers the functionalities of a context-free grammar, allowing the definition of production rules and the analysis of strings.

PushdownAutomaton: This module contains the functionalities of a pushdown automaton, which is an extension of finite automata with an additional stack for temporary storage of information.

These modules are designed to work together, forming a robust and flexible library for the implementation and analysis of FLAT theory models.

\section{OFLAT Platform}
\label{sec:oflat_platform}

The OFLAT platform is an online graphical tool designed for the modeling and analysis of formal language concepts. It provides an intuitive and interactive environment where users can create and manipulate various computational models. The interface consists of a static menu on the left sidebar and a dynamic workspace on the right, which updates based on the selected options.

OFLAT supports the creation and exploration of multiple formal models, including finite automata, pushdown automata, and Turing machines. Users can visually construct these models, define states and transitions, and analyze their properties. The platform offers built-in tools for checking determinism, minimizing models, and identifying structural properties, making it a valuable resource for both learning and research.

By offering a user-friendly graphical interface, OFLAT helps users experiment with formal languages and computation theory concepts in an accessible way. It is designed to be an educational tool that simplifies complex theoretical concepts through visual representation and interactive manipulation.